% Generated draft: requires mathematical and source review.
\documentclass[11pt]{article}
\usepackage[margin=25mm]{geometry}
\usepackage{fontspec}
\setmainfont{Latin Modern Roman}
\usepackage{amsmath,amssymb,mathtools,unicode-math}
\setmathfont{Latin Modern Math}
\usepackage{enumitem,needspace,xurl,hyperref}
\setlength{\parindent}{0pt}
\setlength{\parskip}{5pt}
\setlength{\emergencystretch}{3em}
\newcommand{\N}{\mathbb{N}}
\newcommand{\Z}{\mathbb{Z}}
\newcommand{\Q}{\mathbb{Q}}
\newcommand{\R}{\mathbb{R}}
\newcommand{\C}{\mathbb{C}}
\newcommand{\F}{\mathbb{F}}
\newcommand{\A}{\mathbb{A}}
\newcommand{\PP}{\mathbb{P}}
\newcommand{\E}{\mathbb{E}}
\newcommand{\eps}{\varepsilon}
\newcommand{\dd}{\,\mathrm{d}}
\begin{document}
\section*{UnsolvedMath ID: 1900014}
\textbf{Source problem code:} catalogue entry
\subsection{Definitions and mathematical statement}
For the source's dimension parameter $n\geq3$, the literal task is: prove the existence of a cone realizing a single non-periodic combinatorial structure of a sail.

In a standard contextual model, a pointed cone $C\subset\R^n$ with vertex $0$ gives a lattice convex hull
\[
 K_C=\operatorname{conv}\bigl((C\cap\Z^n)\setminus\{0\}\bigr),
\]
and the boundary faces facing the origin form its sail $\Sigma(C)$. A combinatorial structure would specify the incidence relations of the vertices, edges and higher faces, and realization would require a combinatorial isomorphism from the specified complex to $\Sigma(C)$. The exact source conventions for cones, lattice points and sail boundaries are not supplied, so this is context rather than a recovered definition.

No particular non-periodic complex is supplied, and ``single'' does not specify whether the task is to choose one example with a fully described structure or realize one already prescribed elsewhere in the source. The notion of periodicity, including the acting group and what part of the combinatorial or lattice structure it preserves, is also undefined. Thus the exact existential/constructive assertion cannot be determined from the isolated statement.

The supplied assessment instead discusses realizing \emph{arbitrary} non-periodic combinatorial structures. A universal claim that every admissible non-periodic complex is realizable is stronger than constructing one realization, and is not silently substituted for the literal source. A precise target requires recovering the class of admissible complexes, the periodicity definition, the dimension convention, and which quantifiers are intended.
\subsection{Short English statement}
Construct a cone with the non-periodic sail structure sought by the source, in dimensions $n\geq3$. The original wording asks for a single structure; the later summary broadens this to arbitrary structures without supplying the missing formulation. The summary is explicitly moderate-confidence OPEN-TRIAGE, reports no verified resolution in its search, and mentions related periodic/algebraic realizability work. These are provisional review observations rather than a proof of openness or a complete realization theorem.
\subsection{Sources}
\begin{itemize}
\item UnsolvedMath permanent ID 1900014, source code catalogue entry, ``Geometry of Continued Fractions --- Combinatorial structure of sails.'' Registry curation and metadata: Ulam AI and UnsolvedMath Contributors, CC BY 4.0; updated 2026-08-06.
\item Karpenkov, \emph{Open problems in geometry of continued fractions} (2017), Problem 14: \url{https://arxiv.org/abs/1712.01450}. Supplied source-TeX extraction accessed 2026-07-29. The metadata says the source presents an open item and marks status/release and rights review NEEDS\_REVIEW.
\item The supplied assessment mentions Karpenkov's ``Periodic continued fractions'' and unspecified sail-construction papers without complete dates, titles for the latter, or locators. No further reference or theorem is inferred.
\end{itemize}
\end{document}
